% part: intuitionistic-logic
% chapter: propositions-as-types
% section: types

\documentclass[../../../include/open-logic-section]{subfiles}

\begin{document}

\olfileid{int}{pty}{typ}

\olsection{Tipe}

Terma bukti kaya $(MN)$ dhewe ora nemtokake kanthi tunggal bukti sing
dikode. Iki amarga variabel bebas, yaiku terma bukti saka aksioma,
ora nemtokake !!{formula}s ing aksioma sing cocog. Nanging, yen kita
nemtokake formula kasebut ing konteks~$\Gamma$ lan terma iku tipene
bener, bukti panetepan tipe sing cocog \emph{mesthi} ditemtokake
kanthi tunggal.

\begin{defn}
\emph{Konteks} kanggo terma bukti~$N$ iku himpunan $\Gamma$ sing
masang paling akeh siji formula marang saben variabel, kanthi kanggo
saben variabel bebas saka~$N$, $\Gamma$ ngandhut persis siji pasangan
$x : !A$.
\end{defn}

Gatekna yen dhefinisi ora mbutuhake $\Gamma$ ngandhut $x:!A$
\emph{mung} kanggo variabel bebas saka~$N$. Contone,
$\Gamma = \{x:!A, y:!B, z:!C\}$ iku konteks kanggo $\pair{x}{y}$.

\begin{defn}\ollabel{defn:type}
  Upamane $\Gamma$ iku konteks kanggo~$N$. \emph{Tipe} saka~$N$
  (relatif marang~$\Gamma$) ditetepake kanthi induktif kaya mangkene.
  Jeneng variabel kaiket diganti dhisik supaya anyar tumrap konteks:
  \begin{enumerate}
  \item Tipe $x$ iku $!A$ yen lan mung yen $x:!A \in \Gamma$.
  \item Yen $N$ nduweni tipe~$!B$ relatif marang $x:!A, \Gamma$, tipe $\lambd[\typeof{x}{!A}][N]$ iku $!A \lif !B$.
  \item Yen $N$ nduweni tipe $!A \lif !B$ lan $M$ nduweni tipe~$!A$,
    $(NM)$ nduweni tipe~$!B$
  \item Yen $N$ nduweni tipe $!A$ lan $M$ nduweni tipe~$!B$,
    $\pair{N}{M}$ nduweni tipe $!A \land !B$.
  \item Yen $N$ nduweni tipe~$!A_1 \land !A_2$, $\proj{i}{N}$ nduweni tipe~$!A_i$.
  \item Yen $N$ nduweni tipe $!A$,
    $\inj[!B]{i}{N}$ nduweni tipe $!A \lor !B$ yen $i=1$ lan tipe
    $!B \lor !A$ yen $i=2$.
  \item Yen $M$ nduweni tipe $!A \lor !B$, lan $N_1$ nduweni tipe
    $!C$ relatif marang $x_1:!A, \Gamma$ lan $N_2$ nduweni tipe~$!C$
    relatif marang $x_2:!B, \Gamma$, $\dcase{M}{x_1}{N_1}{x_2}{N_2}$
    nduweni tipe~$!C$.
  \item Yen $N$ nduweni tipe~$\lfalse$,
    $\abort{!A}{N}$ nduweni tipe~$!A$.
  \end{enumerate}
\end{defn}

\begin{prop}
Saben terma bukti~$N$ nduweni paling akeh siji tipe relatif marang
konteks~$\Gamma$ kanggo~$N$. Saben terma sing tipene bener nduweni
persis siji tipe.
\end{prop}

\begin{proof}
  Kanthi induksi marang~$N$.
\end{proof}

Kita nulis $\Gamma \Sequent N : !A$ yen tipe $N$ relatif marang
$\Gamma$ iku~$!A$. Klausa ing \olref{defn:type} mesthine wis katon
lumrah: klausa kasebut persis aturan \Log{tN2i} kanthi siji bedane
cilik, yaiku konteks~$\Gamma$ padha ing sakabehe. Kita bisa ngrumusake
dhefinisi kasebut minangka kalkulus \Log{tN3i} kanthi aturan ing
\olref{tab:tN3ip}.

\subfile{rules-tN3}

Terma bukti iku terma ing sawijining versi kalkulus $\lambd$, yaiku
\emph{kalkulus $\lambd$ mawa tipe kanthi prodhuk, gunggung, lan tipe kosong}.
Kalkulus $\lambd$ klasik mung nduweni terma sing digawe saka variabel,
\emph{abstraksi $\lambd$}~$\lambd[x][N]$ lan aplikasi $(MN)$,
lan ora nduweni anotasi~$!A$ ing abstraksi $\lambd$ mawa tipe
$\lambd[\typeof{x}{!A}][N]$. Kalkulus $\lambd$ iku modhel komputasi
sing lengkap Turing (yaiku saben fungsi komputabel parsial bisa
diwakili dening terma ing kalkulus kasebut). Versi kalkulus $\lambd$
sing diwenehake dening terma bukti lan aturan tipe iku modhel
komputasi sing diwatesi, nanging gagasan dhasare padha. Wates kasebut
ditetepake dening panetepan tipe: mung terma \emph{sing tipene bener},
yaiku terma bukti mawa tipe (relatif marang konteks~$\Gamma$), sing
dadi terma sah. Contone, $\lambd[x][(xx)]$ iku terma sah ing kalkulus
$\lambd$ klasik tanpa tipe, nanging $\lambd[\typeof{x}{!A}][(xx)]$
ora bisa nduweni tipe sing bener kanggo $!A$ endi wae. Iki amarga
aturan tipe mbutuhake, supaya $(xx)$ nduweni tipe, $x$ kapisan kudu
mawa tipe $!A \lif !B$ lan $x$ kapindho mawa tipe~$!A$, lan iki
ora mungkin amarga $!A$ iku sub-!!{formula} sejati saka~$!A \lif !B$.

Kanthi intuisi, kita bisa nganggep tipe minangka jinis barang sing
dijenengi dening terma mawa tipe kasebut. Mula terma mawa tipe
$!A \land !B$ nemtokake pasangan kanthi komponen kapisan mawa
tipe~$!A$ lan komponen kapindho mawa tipe~$!B$, yaiku !!a{element}
saka prodhuk~$!A \times !B$. Terma mawa tipe $!A \lif !B$ iku
fungsi saka barang mawa tipe $!A$ menyang barang mawa tipe~$!B$.
Terma mawa tipe $!A \lor !B$ iku barang mawa tipe $!A$ utawa barang
mawa tipe~$!B$, bebarengan karo informasi endi sing dipilih.
Iki kaya gabungan pisah saka rong himpunan $!A$ lan $!B$ sing
ditetepake minangka $\{\tuple{i,x} : i = 1 \text{ utawa } 2,
x \in !A \text{ yen } i = 1, x \in !B \text{ yen } i=2\}$.
Tipe $\lfalse$ iku himpunan kosong, yaiku yen sawijining terma mawa
tipe kasebut, terma iku ora bisa menehi jeneng barang apa wae.
Amarga dedhuksi alami (kalebu \Log{tN3i}) konsisten, ora ana
!!{proof} saka~$\lfalse$ kajaba ana asumsi mbukak, mula ora ana
terma bukti katutup (terma tanpa variabel bebas) mawa tipe~$\lfalse$.

Operator kalkulus $\lambd$ mawa tipe ngasilake terma sing, kanthi
intuisi, menehi jeneng barang saka jinis sing bener, yen terma sing
diwenehi operator kasebut menehi jeneng barang tartamtu. Contone,
``yen $N_1 : !A$ lan $N_2 : !B$, mula $\pair{N_1}{N_2} : !A \land !B$''
nyatakake yen nalika $N_1$ menehi jeneng barang mawa tipe~$!A$ lan
$N_2$ menehi jeneng barang mawa tipe~$!B$, $\pair{N_1}{N_2}$ menehi
jeneng barang mawa tipe $!A \land !B$.

\end{document}
